% GENERATED FILE. Edit canonical lessons, then run npm run generate:books.
% canonical-source-hash: fnv1a64:0af8b534511f7ca8
% canonical-lessons: KA-C446-kaludari, KA-C446-oladari, KA-C446-suttudari, KA-C446-dari-tappu, KA-C446-hejje

\chapter{Shortcuts and Detours}
\label{ch:ka-shortcuts-and-detours}
\hlchaptermodality{446}

\begin{chapteropening}
\textbf{By the end of this chapter:} \emph{I can say a footpath, a shortcut, a detour, to lose one's way and a footstep.}

Complete the last lesson of chapter 446: I can say a footpath, a shortcut, a detour, to lose one's way and a footstep.
\end{chapteropening}

\section[\texorpdfstring{\kn{ಕಾಲುದಾರಿ} (kāludāri)}{ಕಾಲುದಾರಿ (kāludāri)}]{\kn{ಕಾಲುದಾರಿ} (kāludāri) — a footpath}

\label{lesson:KA-C446-kaludari}



\begin{quote}
Before the new one: say the Kannada for to move, then the Kannada for to approach.
\end{quote}

\subsection*{You'll want to know: \kn{ಕಾಲುದಾರಿ}}
\textbf{\kn{ಕಾಲುದಾರಿ}} — \emph{kāludāri} — \textquotedblleft{}a footpath\textquotedblright{}.

\begin{quote}
\kn{ಹೊಲಗಳ} \kn{ನಡುವೆ} \kn{ಒಂದು} \kn{ಕಾಲುದಾರಿ} \kn{ಇದೆ}.
\end{quote}

\begin{quote}
\emph{holagaḷa naḍuve ondu kāludāri ide.}
\end{quote}

\begin{quote}
There is a footpath between the fields.
\end{quote}

\kn{ಕಾಲು}, foot, plus \kn{ದಾರಿ}, way: a way for feet only.

\begin{cousinweb}
\textbf{\kn{ಕಾಲುದಾರಿ}} (\emph{kāludāri}) is a native Kannada compound of \emph{kāl}, \textquotedblleft{}foot, leg\textquotedblright{} (DEDR 1479), and \kn{ದಾರಿ}, \textquotedblleft{}a way\textquotedblright{} (DEDR 3170).
\end{cousinweb}

\begin{grammarlens}[title={What you've built}]
One more word for getting around.
\end{grammarlens}

\subsection*{Your turn}
\begin{itemize}
\raggedright
  \item \emph{Say it:} \emph{kāludāri}
  \item \emph{Say it:} \emph{kāludāri}, once more
  \item \emph{Say it:} \emph{holagaḷa naḍuve ondu kāludāri ide.} — There is a footpath between the fields.
  \item \emph{Recall:} say the Kannada for to move, then the Kannada for to approach, then say \emph{kāludāri} again
\end{itemize}

\begin{tcolorbox}[breakable,colback=teal!4,colframe=teal!35!black,title={Before you move on}]
What does \kn{ಕಾಲುದಾರಿ} mean? (\textquotedblleft{}A footpath\textquotedblright{}.) Say it once more.
\end{tcolorbox}
\section[\texorpdfstring{\kn{ಒಳದಾರಿ} (oḷadāri)}{ಒಳದಾರಿ (oḷadāri)}]{\kn{ಒಳದಾರಿ} (oḷadāri) — a shortcut; a back way}

\label{lesson:KA-C446-oladari}



\begin{quote}
Before the new one: say the Kannada for to approach, then the Kannada for a footpath.
\end{quote}

\subsection*{You'll want to know: \kn{ಒಳದಾರಿ}}
\textbf{\kn{ಒಳದಾರಿ}} — \emph{oḷadāri} — \textquotedblleft{}a shortcut; a back way\textquotedblright{}.

\begin{quote}
\kn{ಒಳದಾರಿಯಲ್ಲಿ} \kn{ಹೋದರೆ} \kn{ಹತ್ತು} \kn{ನಿಮಿಷ} \kn{ಉಳಿಯುತ್ತದೆ}.
\end{quote}

\begin{quote}
\emph{oḷadāriyalli hōdare hattu nimiṣa uḷiyuttade.}
\end{quote}

\begin{quote}
If you take the shortcut, you save ten minutes.
\end{quote}

\kn{ಒಳ}, inner, plus \kn{ದಾರಿ}, way: a way through the inside.

\begin{cousinweb}
\textbf{\kn{ಒಳದಾರಿ}} (\emph{oḷadāri}) is a native Kannada compound of \emph{oḷa}, \textquotedblleft{}the inside, inner\textquotedblright{} (DEDR 698), and \kn{ದಾರಿ}, \textquotedblleft{}a way\textquotedblright{} (DEDR 3170).
\end{cousinweb}

\begin{grammarlens}[title={What you've built}]
One more word for getting around.
\end{grammarlens}

\subsection*{Your turn}
\begin{itemize}
\raggedright
  \item \emph{Say it:} \emph{oḷadāri}
  \item \emph{Say it:} \emph{oḷadāri}, once more
  \item \emph{Say it:} \emph{oḷadāriyalli hōdare hattu nimiṣa uḷiyuttade.} — If you take the shortcut, you save ten minutes.
  \item \emph{Recall:} say the Kannada for to approach, then the Kannada for a footpath, then say \emph{oḷadāri} again
\end{itemize}

\begin{tcolorbox}[breakable,colback=teal!4,colframe=teal!35!black,title={Before you move on}]
What does \kn{ಒಳದಾರಿ} mean? (\textquotedblleft{}A shortcut; a back way\textquotedblright{}.) Say it once more.
\end{tcolorbox}
\section[\texorpdfstring{\kn{ಸುತ್ತುದಾರಿ} (suttudāri)}{ಸುತ್ತುದಾರಿ (suttudāri)}]{\kn{ಸುತ್ತುದಾರಿ} (suttudāri) — a detour; a roundabout way}

\label{lesson:KA-C446-suttudari}



\begin{quote}
Before the new one: say the Kannada for a footpath, then the Kannada for a shortcut.
\end{quote}

\subsection*{You'll want to know: \kn{ಸುತ್ತುದಾರಿ}}
\textbf{\kn{ಸುತ್ತುದಾರಿ}} — \emph{suttudāri} — \textquotedblleft{}a detour; a roundabout way\textquotedblright{}.

\begin{quote}
\kn{ಸೇತುವೆ} \kn{ಮುಚ್ಚಿದ್ದರಿಂದ} \kn{ನಾವು} \kn{ಸುತ್ತುದಾರಿಯಲ್ಲಿ} \kn{ಹೋದೆವು}.
\end{quote}

\begin{quote}
\emph{sētuve mucciddarinda nāvu suttudāriyalli hōdevu.}
\end{quote}

\begin{quote}
Because the bridge was closed we went the long way round.
\end{quote}

\kn{ಸುತ್ತು}, a round, plus \kn{ದಾರಿ}, way.

\begin{cousinweb}
\textbf{\kn{ಸುತ್ತುದಾರಿ}} (\emph{suttudāri}) is a native Kannada compound of \emph{suttu}, \textquotedblleft{}to go round; a turn\textquotedblright{} (DEDR 2715), and \kn{ದಾರಿ}, \textquotedblleft{}a way\textquotedblright{} (DEDR 3170).
\end{cousinweb}

\begin{grammarlens}[title={What you've built}]
One more word for getting around.
\end{grammarlens}

\subsection*{Your turn}
\begin{itemize}
\raggedright
  \item \emph{Say it:} \emph{suttudāri}
  \item \emph{Say it:} \emph{suttudāri}, once more
  \item \emph{Say it:} \emph{sētuve mucciddarinda nāvu suttudāriyalli hōdevu.} — Because the bridge was closed we went the long way round.
  \item \emph{Recall:} say the Kannada for a footpath, then the Kannada for a shortcut, then say \emph{suttudāri} again
\end{itemize}

\begin{tcolorbox}[breakable,colback=teal!4,colframe=teal!35!black,title={Before you move on}]
What does \kn{ಸುತ್ತುದಾರಿ} mean? (\textquotedblleft{}A detour; a roundabout way\textquotedblright{}.) Say it once more.
\end{tcolorbox}
\section[\texorpdfstring{\kn{ದಾರಿ} \kn{ತಪ್ಪು} (dāri tappu)}{ದಾರಿ ತಪ್ಪು (dāri tappu)}]{\kn{ದಾರಿ} \kn{ತಪ್ಪು} (dāri tappu) — to lose one's way}

\label{lesson:KA-C446-dari-tappu}



\begin{quote}
Before the new one: say the Kannada for a shortcut, then the Kannada for a detour.
\end{quote}

\subsection*{You'll want to know: \kn{ದಾರಿ} \kn{ತಪ್ಪು}}
\textbf{\kn{ದಾರಿ} \kn{ತಪ್ಪು}} — \emph{dāri tappu} — \textquotedblleft{}to lose one's way\textquotedblright{}.

\begin{quote}
\kn{ಕತ್ತಲೆಯಲ್ಲಿ} \kn{ನಾವು} \kn{ಕಾಡಿನಲ್ಲಿ} \kn{ದಾರಿ} \kn{ತಪ್ಪಿದೆವು}.
\end{quote}

\begin{quote}
\emph{kattaleyalli nāvu kāḍinalli dāri tappidevu.}
\end{quote}

\begin{quote}
In the dark we lost our way in the forest.
\end{quote}

Past (I …): \kn{ದಾರಿ} \kn{ತಪ್ಪಿದೆ}, \emph{dāri tappide}.

Word for word the way slips: \kn{ತಪ್ಪು} is to slip or to go wrong.

\begin{cousinweb}
\textbf{\kn{ದಾರಿ} \kn{ತಪ್ಪು}} (\emph{dāri tappu}) joins two native Dravidian words: \kn{ದಾರಿ}, \textquotedblleft{}a way\textquotedblright{} (DEDR 3170), and \kn{ತಪ್ಪು}, \textquotedblleft{}to slip, to miss, to go wrong\textquotedblright{} (DEDR 3071).
\end{cousinweb}

\begin{grammarlens}[title={What you've built}]
One more word for getting around.
\end{grammarlens}

\subsection*{Your turn}
\begin{itemize}
\raggedright
  \item \emph{Say it:} \emph{dāri tappu}
  \item \emph{Say it:} \emph{dāri tappu}, once more
  \item \emph{Say it:} \emph{kattaleyalli nāvu kāḍinalli dāri tappidevu.} — In the dark we lost our way in the forest.
  \item \emph{Recall:} say the Kannada for a shortcut, then the Kannada for a detour, then say \emph{dāri tappu} again
\end{itemize}

\begin{tcolorbox}[breakable,colback=teal!4,colframe=teal!35!black,title={Before you move on}]
What does \kn{ದಾರಿ} \kn{ತಪ್ಪು} mean? (\textquotedblleft{}To lose one's way\textquotedblright{}.) Say it once more.
\end{tcolorbox}
\section[\texorpdfstring{\kn{ಹೆಜ್ಜೆ} (hejje)}{ಹೆಜ್ಜೆ (hejje)}]{\kn{ಹೆಜ್ಜೆ} (hejje) — a footstep; a step}

\label{lesson:KA-C446-hejje}



\begin{quote}
Before the new one: say the Kannada for a detour, then the Kannada for to lose one's way.
\end{quote}

\subsection*{You'll want to know: \kn{ಹೆಜ್ಜೆ}}
\textbf{\kn{ಹೆಜ್ಜೆ}} — \emph{hejje} — \textquotedblleft{}a footstep; a step\textquotedblright{}.

\begin{quote}
\kn{ಮರಳಿನ} \kn{ಮೇಲೆ} \kn{ಆನೆಯ} \kn{ಹೆಜ್ಜೆಗಳು} \kn{ಕಂಡವು}.
\end{quote}

\begin{quote}
\emph{maraḷina mēle āneya hejjegaḷu kaṇḍavu.}
\end{quote}

\begin{quote}
Elephant footprints could be seen on the sand.
\end{quote}

Plural: \kn{ಹೆಜ್ಜೆಗಳು}, \emph{hejjegaḷu}.

\kn{ಹೆಜ್ಜೆ} \kn{ಹಾಕು}, to take steps, to walk along.

\begin{cousinweb}
\textbf{\kn{ಹೆಜ್ಜೆ}} (\emph{hejje}) \textquotedblleft{}a footprint, a step\textquotedblright{} is listed, beside older \emph{pajje}, in the appendix of Indo-Aryan loans in DEDR (DEDR App. 6), which compares Sanskrit \emph{padyā}, \textquotedblleft{}footsteps\textquotedblright{} (CDIAL 7778); the exact path is not settled.
\end{cousinweb}

\begin{grammarlens}[title={What you've built}]
One more word for getting around.
\end{grammarlens}

\subsection*{Your turn}
\begin{itemize}
\raggedright
  \item \emph{Say it:} \emph{hejje}
  \item \emph{Say it:} \emph{hejje}, once more
  \item \emph{Say it:} \emph{maraḷina mēle āneya hejjegaḷu kaṇḍavu.} — Elephant footprints could be seen on the sand.
  \item \emph{Recall:} say the Kannada for a detour, then the Kannada for to lose one's way, then say \emph{hejje} again
\end{itemize}

\begin{tcolorbox}[breakable,colback=teal!4,colframe=teal!35!black,title={Before you move on}]
What does \kn{ಹೆಜ್ಜೆ} mean? (\textquotedblleft{}A footstep; a step\textquotedblright{}.) Say it once more.
\end{tcolorbox}
